% ============================================================
% Aarushi's report sections (Q0.1, Q1.1--Q1.4, Q2.1--Q2.3)
% Paste each block into the matching response box of the template.
% Tables use plain tabular + \hline (no extra packages needed).
% If the template does not already load graphicx, add \usepackage{graphicx}.
% ============================================================

% ---------------- Q0.1 ----------------
\textbf{Aarushi Singh (as469):} Set up the repository and shared pipeline (data
loading, fixed stratified validation split, evaluation and submission scripts).
Ran the model comparison (logistic regression, linear SVM, Naive Bayes, gradient
boosting on stylometric features) and hyperparameter searches. Discovered and
analyzed the blank-line formatting artifact. Ran robustness tests (typos,
homoglyphs, word dropping, newline removal) and error analysis. Ran the
pseudo-labeling experiments. Fine-tuned DistilRoBERTa, built the GPT-2 perplexity
and Binoculars features, built the rank-averaged ensembles, and managed the
Kaggle submissions.

% ---------------- Q1.1 ----------------
We compared four representations. \textbf{(1) Sparse $n$-grams:} scikit-learn's
default tokenizer with lowercasing and no stemming or stopword removal, since
function-word frequency proved informative. Word 1--2-grams
(\texttt{min\_df}=2, $\le$200k features) capture phrasing such as ``overall'' or
``such as''; character 3--5-grams (\texttt{char\_wb}, \texttt{min\_df}=3,
$\le$300k features) capture sub-word patterns. Both use sublinear TF-IDF.
\textbf{(2) Stylometric features:} 26 statistics (word/sentence length,
type--token ratio, function-word, punctuation and newline rates).
\textbf{(3) Fine-tuned DistilRoBERTa} with its byte-level BPE tokenizer,
truncated to 256 tokens. \textbf{(4) Perplexity features:} 13 statistics of
GPT-2's per-token log-probabilities, entropy and rank, plus a Binoculars-style
ratio. After finding that blank lines leaked the label (0\% of human vs.\ 66\%
of AI training passages) and did not transfer to the test set, we
\textbf{collapsed all whitespace} before the transformer and perplexity models,
so they cannot use formatting.

% ---------------- Q1.2 ----------------
\begin{center}
\footnotesize
\begin{tabular}{|p{3.0cm}|p{4.2cm}|c|c|p{2.5cm}|}
\hline
\textbf{Model} & \textbf{Text representation} & \textbf{Train AUROC} & \textbf{Val.\ AUROC} & \textbf{Train time} \\
\hline
Logistic regression & Word 1--2g + char 3--5g TF-IDF & 1.000 & 0.9506 & 9.7 s \\
Logistic regression & Word 1--2g TF-IDF & 1.000 & 0.9413 & 2.9 s \\
Logistic regression & Char 3--5g TF-IDF & 0.998 & 0.9386 & 10.1 s \\
Linear SVM (calibrated) & Word + char TF-IDF & 1.000 & 0.9504 & 8.1 s \\
Multinomial Naive Bayes & Word 1--2g counts & 0.999 & 0.8682 & 1.6 s \\
Gradient boosting & 26 stylometric features & 1.000 & 0.9945$^{*}$ & 6.4 s \\
Gradient boosting & 13 GPT-2 perplexity features & 0.9995 & 0.9706 & $\sim$9 min feat.\ + $<$10 s fit \\
DistilRoBERTa (2 epochs) & Subword tokens, whitespace collapsed & 0.9951 & 0.9867 & 10.9 min (M4 GPU) \\
\textbf{Ensemble (0.5/0.3/0.2)} & DistilRoBERTa + perplexity GBM + TF-IDF LR & -- & \textbf{0.9940} & sum of parts \\
\hline
\end{tabular}
\end{center}
{\footnotesize $^{*}$Relies on the newline artifact: 0.9418 without the newline
feature, and only 0.858 on the public leaderboard.}

% ---------------- Q1.3 ----------------
\begin{center}
\small
\begin{tabular}{|l|p{8.2cm}|l|}
\hline
\textbf{Model} & \textbf{Parameters / values searched} & \textbf{Chosen setting} \\
\hline
TF-IDF logistic regression & $C \in \{0.25, 0.5, 1, 2, 4, 8, 16, 32\}$ & $C = 4$ \\
Stylometric gradient boosting & \texttt{learning\_rate} $\in \{0.03, 0.05, 0.1\}$ $\times$
\texttt{max\_iter} $\in \{100, 300, 600\}$ $\times$ \texttt{max\_leaf\_nodes} $\in \{15, 31, 63\}$
(27 settings) & 0.05 / 300 / 31 \\
Ensemble weights & Transformer / perplexity / TF-IDF weights on a 0.1 grid summing to 1 & 0.5 / 0.3 / 0.2 \\
\hline
\end{tabular}
\end{center}

Every setting was trained on the same 8,000-row training split and scored on the
fixed, stratified 2,000-row validation split (seed 42); Kaggle scores were never
used for tuning. For logistic regression, validation AUROC rose from 0.932
($C=0.25$) to a plateau of 0.949--0.951 for $C \ge 2$; we kept $C=4$, since
$C=8$ was only 0.0006 higher, which is within noise. The GBM grid ranged from
0.986 to 0.995; we kept the default (0.9945), because the best setting (0.9953)
differed by under 0.001 and used twice as many trees. Ensemble weights were
picked from a flat region: the top 10 combinations all scored 0.993--0.994, so
the choice is not fragile.

% ---------------- Q1.4 ----------------
Our final submission is a \textbf{rank-average ensemble of fine-tuned
DistilRoBERTa (0.5), a GPT-2 perplexity GBM (0.3) and word+char TF-IDF logistic
regression (0.2)}, with 0.9940 validation AUROC and \textbf{0.921 on the public
leaderboard}. Validation score alone was not enough: the stylometric GBM had the
best validation AUROC (0.9945), but its strength came from a blank-line artifact,
and it scored 0.858 on Kaggle, the largest validation-to-test drop of any model
($-0.137$). We therefore favored models that see no formatting and whose errors
complement each other. DistilRoBERTa alone gained 3 points on Kaggle over every
classical model (0.894). Adding TF-IDF raised it to 0.902, and adding perplexity
features, which do not depend on topic or formatting, raised it to 0.921.
Training costs about 25 minutes on a laptop GPU. The top weightings were stable
($\pm$0.001), suggesting the ensemble is not tuned to validation noise.

% ---------------- Q2.1 ----------------
\begin{center}
\begin{tabular}{|r|c|c|l|}
\hline
\textbf{Training rows} & \textbf{Training AUROC} & \textbf{Validation AUROC} & \textbf{Notes} \\
\hline
500   & 1.000 & 0.870 $\pm$ 0.007 & mean of 3 seeds \\
1,000 & 1.000 & 0.897 $\pm$ 0.003 & mean of 3 seeds \\
2,000 & 1.000 & 0.924 $\pm$ 0.001 & mean of 3 seeds \\
5,000 & 1.000 & 0.943 $\pm$ 0.001 & mean of 3 seeds \\
8,000 & 1.000 & 0.951 & single fit (full split) \\
\hline
\end{tabular}
\end{center}
{\footnotesize Model: word+char TF-IDF logistic regression, $C=4$, fixed
hyperparameters and the same 2,000-row validation set throughout.}

% ---------------- Q2.2 ----------------
\begin{center}
\includegraphics[width=0.85\linewidth]{learning_curve_train_val.png}
\end{center}
% (copy results/learning_curve_train_val.png next to your .tex file, or adjust the path)

% ---------------- Q2.3 ----------------
Training AUROC is 1.000 at every size: with up to 500k sparse features and at
most 8,000 passages, the logistic regression can always separate its training
set perfectly, so the model is clearly \textbf{not underfitting}. The gap between
training and validation AUROC indicates \textbf{overfitting (high variance)}, but
it shrinks steadily as data grows, from 0.130 at 500 rows to 0.049 at 8,000.
Validation AUROC is still rising at 8,000 rows ($+0.008$ from 5k to 8k), roughly
linearly in $\log n$, so \textbf{more labeled data would still help}, though with
diminishing returns per example. Stronger regularization does not help: lowering
$C$ reduced both training and validation AUROC (0.932 at $C=0.25$), suggesting
the model is limited by data rather than capacity, which also motivated our
pseudo-labeling experiment. Seed-to-seed variability also falls with size
($\pm$0.007 to $\pm$0.001).
